\documentclass[11pt]{article}
\usepackage[margin=1in]{geometry}
\usepackage{microtype}
\usepackage{amsmath}
\usepackage[numbers,sort&compress]{natbib}
\usepackage{hyperref}
\usepackage{parskip}

\bibliographystyle{unsrtnat}

\begin{document}

\section*{Background and Central Hypothesis}

In the following we present a revised proposal which builds and expands on our previous proposal. We further elaborate on the two steps/research questions we presented in the previous proposal, present a deeper theoretical motivation and develop concrete experimental setup to test the refined research questions. 

We conceptualize intelligent agents as solving problems through a sequence of two steps. When faced with a concrete problem (like AlphaZero with a board state, or a language model with a math problem), the agent has an initial set of priors that inform the next action it takes (the next move to be played or the next token to be generated). The agent then uses an exploratory phase to accumulate information and shifts its prior distribution over the next action towards something (hopefully) better. Formally, this is the structure of Bayes-adaptive sequential decision-making \citep{duff2002,ghavamzadeh2015}.

In AlphaZero \citep{silver2018alphazero}, this exploratory phase is hardcoded through a Monte-Carlo tree search of the game tree where each position is scored using the learned value function. Once this exploration is finished, there is also a hardcoded algorithm to use this valued game tree to improve its initial prior distribution over the next move.

In language models with ``reasoning'' \citep{deepseek2025r1}, this exploratory phase is simply the chain-of-thought the model uses, and is determined purely by the model itself. Once this phase is completed, the model leverages its in-context learning (ICL) abilities \citep{brown2020gpt3,xie2022icl} to improve its answer. The key difference from AlphaZero is that both how the exploration is executed, and how the exploration is used to improve the prior distribution over actions, is learned during training, instead of being a hardcoded algorithm. This difference informs our central hypothesis.

\paragraph{Central hypothesis.} When training jointly shapes both \emph{(i)} how an agent explores and \emph{(ii)} how it updates its policy from the exploration --- as in RL-trained reasoning models, but not in AlphaZero --- gradient pressure incentivises two internal representations to co-develop: an explicit representation of the agent's beliefs over task-relevant latents, and a causal model of the environment in which the agent itself figures as a distinct, well-separated cause. Put simply, we expect  that the pressure to explore optimally and to learn from that exploration well will lead to the emergence of an internal belief representation and a causal self representation that is separate from the world representation, and that the self and world representations factor the dynamics of the activations of the agent.

\paragraph{Why we expect this.} Optimal exploration is exploration that maximally increases information about the agent's task-relevant beliefs \citep{russo2014ids}. To do this, the agent must choose exploratory actions informed by its current beliefs, which motivates an internal representation of those beliefs and circuits that consult them to guide exploration. The second representation is required for a related reason: when the agent's own actions are part of what produces its next observation (as is the case for every token of a chain-of-thought), correctly attributing the outcome to its real causes requires the agent to distinguish what it produced itself from what is due to the world. Without this distinction, gradient updates conflate self-generated patterns with genuine signal about the task. When both the exploration and the update from it are themselves learned, gradient descent has every reason to develop both representations.

The first half of this hypothesis --- the emergence of internal belief representations under exploration pressure --- is grounded in two research communities, the ML and developmental cognitive science. On the ML side, in the setting of sequence prediction in the context of Hidden Markov Models (HMMs), transformers trained as next-token predictors develop linearly recoverable internal representations of the distribution over the hidden states \citep{shai2024belief}. In active settings, RL-trained recurrent networks develop similar representations over hidden states \citep{hennig2023emergence}, and on behavioural and representational grounds, meta-trained agents converge to Bayes-optimal solutions whose state must encode a posterior over the task \citep{mikulik2020bayesopt,ortega2019meta}.

Developmental cog sci provides complementary findings: infants attend preferentially to stimuli at intermediate complexity \citep[the Goldilocks effect,][]{kidd2012goldilocks}, select informative causal interventions \citep{lapidow2020informative}, track their own learning progress \citep{poli2020infants}, adapt their search to the ecological structure of the environment \citep{ruggeri2022ecological}, and show pupillometric signatures of prediction error minimisation in response to stimuli \citep{zhang2019prediction}. What this work cannot resolve is the internal computations underlying these behaviours: infants behave as if they are optimising for information, but the objectives themselves and the internal states that implement them are not directly observable from outside.

The second half of the hypothesis --- the agent's representation of itself as a distinct cause --- has received much less direct attention in the ML literature, though circumstantial evidence is accumulating. Post-trained language models exhibit behavioural self-awareness \citep{betley2025tellme}, and self-recognition that correlates with self-preference \citep{panickssery2024self}. The forthcoming work of one of us (Asvin) with Jack Lindsey finds that post-trained models can detect when they are reading their own on-policy generations and modulate behaviour accordingly. What none of this work does is test, mechanistically and systematically, whether a self representation in the sense the hypothesis requires actually exists, factorises from belief representations, and is causally load-bearing in the forward pass.

Developmental cognitive science has a longer-standing engagement with the self / world question, through a mechanism that maps especially cleanly onto the one we appeal to in language models. For example,  within hours of birth, infants alter their sucking rhythm to produce preferred sounds \citep{decasper2009lateralized}. Newborns will also raise an arm into a light beam and attend to the resulting hand shadow \citep{vandermeer1997limelight}. By three months, infants kick more vigorously when their ankle is tethered to a mobile that responds to their movement than under a yoked non-contingent control \citep{rovee1969conjugate}, and generalise the same principle across modalities and effectors \citep{lee2013contingent}. The mechanism running through all of these is temporal contingency detection: the infant learns to identify its own actions by their reliable, short-latency coupling to reafferent feedback. This is precisely the causal-sampling property we use as the signature of self in our toy and LLM analyses --- a self-emitted token is recognised as self by being a sample from the model's own prior output distribution, in the same way that an infant's movement is recognised as self by its temporal coupling to the motor intention that preceded it. The adult signature of the same computation is sensory attenuation: self-produced stimulation is perceptually damped relative to externally produced stimulation \citep{blakemore1998tickle}. In our terms this is an actor-gated belief update in which self-caused input is discounted as evidence, and the actor-gated likelihood factor introduced in the analysis section makes this prediction quantitative.

More striking still is the prenatal evidence, an argument developed at greater length by one of us \citep{rutar2026prenatal}. From the onset of foetal movement around 7.5 weeks gestation \citep{prechtl1985ultrasound}, the structure of motor output increases rapidly: early vermicular movements
(around 7.5 weeks GA; \citep{prechtl1985ultrasound}) give way to increasingly
fractionated, spatially selective limb and joint movements (such as general movements,
startles, and brief myoclonic limb twitches; \citep{devries1982emergence,
piontelli2010conclusions, tiriac2014selfgenerated, lindemann2016spindle}) by 20 weeks
\citep{devries1982emergence}, consistent with a shift from coarse, metabolically cheap sampling toward more targeted exploration of the foetus's own sensorimotor space. The earliest behavioural evidence of a primitive self-model appears at approximately the same time as evidenced by the following and similar findings: by 22 weeks foetal hand trajectories toward the mouth follow more direct paths than those toward the eyes \citep{zoia2007evidence}, implying a mapping from motor commands to their expected sensory consequences on a self-structured target. The co-emergence of structured exploration and a primitive causal model of the body is consistent with the pattern our hypothesis predicts, though, crucially, as a correlational observation it cannot show that one drives the other. Whether the structure of exploration itself drove the formation of the self-model, or whether both reflect a common maturational programme, cannot be determined from developmental data alone.

The contribution of the proposed work --- spanning both halves of the hypothesis --- is to do what neither literature can do alone. On the developmental side, no body of behavioural evidence, however rich, can confirm whether the optimisation objectives it imputes to learning systems are the ones actually being optimised, or whether the internal representations underlying those objectives have the structure we hypothesise: developmental data cannot vary the training regime, intervene on representations, or access the ground truth. On the ML side, despite the circumstantial evidence cited above, no prior work has systematically tested whether self representations in the sense we define them --- actor-identity representations arising from the causal-sampling structure of self-produced tokens, factorised from belief and causally load-bearing in the forward pass --- exist or co-develop with belief representations under the right training pressure. We address both gaps by using circuit-level probes and interventions across systematically varied training regimes to test whether the belief representations \citep{shai2024belief,hennig2023emergence} are entangled with the self representations \citep{betley2025tellme}, whether the two factorise as a unified Self/World structure, and whether this structure emerges from the joint training pressure on exploration and update that distinguishes RL-trained reasoning models from AlphaZero-style systems. We pursue this first in a controlled toy setting where the analytical ground truth is computable, and then extend the same probe suite to open-source language models trained with environment interaction. 

In short, developmental research and ML research inform and constrain one another in a mutually reinforcing way: developmental research provides the empirical patterns and conceptual vocabulary that motivate our hypothesis, and our work in turn aims to clarify which causal mechanisms can produce those patterns and which cannot, sharpening the hypotheses the developmental community is itself testing.

\section*{Experimental Design for the First Half of the Year}

In order to isolate the various factors above cleanly, we propose a first experiment where we modify the setup in \citep{shai2024belief}. They train a transformer network on a next-token prediction task, where the tokens are generated by a Hidden Markov Model (HMM). They find that purely through being trained on the next prediction task, the model learns to internally represent a distribution over the hidden states of the HMM, and to learn correct update rules over this distribution from one token to the next. This is the simplest model of a \emph{pure predictor}. The model cannot influence the token stream in any way, and is constrained to passive prediction. As such, there is no "self" or "agency" to be learned. 

We propose to modify the set-up minimally in order to introduce agency. We maintain the same HMM model as before for the world, but introduce an ability for the model to "act". The token stream is now of the form $(a_1,o_1,\dots,a_n,o_n)$ where the $a_i$ are "actions", i.e., they are tokens sampled from the prediction of the neural network on the task. These actions \emph{cause} a transition in the hidden state of the HMM, and optionally provide a reward to the agent. The $o_i$ are observations as before and emitted when the HMM undergoes a state transition.

The generative process underlying this token stream factorises naturally into two coupled Markov processes that run in parallel. A \emph{World} chain governs the hidden environment state $w_t$, with transitions driven by the agent's actions ($w_{t+1} \sim P(\cdot \mid w_t, a_t)$) and emissions giving the observations $o_i$. A \emph{Self} chain governs the agent's own policy state $s_t$, with transitions driven by both the action it emits and the observation it subsequently receives ($s_{t+1} = g(s_t, a_t, o_{t+1})$) and emissions giving the actions $a_i$. At every time step both chains have a state, and they are coupled only through the visible tokens they exchange --- the actions and observations act as a Markov blanket between Self and World. The token-position parity simply reflects this exchange: odd positions carry actions emitted by Self, even positions carry observations emitted by World. Structures of this kind appear in the literature under different names (the World half is an input-output HMM \citep{bengio1995iohmm} with actions as inputs, and the joint system is a coupled HMM \citep{brand1996coupled}) though the explicit Self / World decomposition is not standard. 

We refer to this joint generative model as the \textbf{SWITCH} (Self/World Interactive Turn-taking Coupled HMM). We make it explicit here because the central hypothesis predicts that the trained transformer should recover this factorisation \emph{internally}: a Self subspace whose activations encode the agent's policy state, a separable World subspace whose activations encode the Bayes-adaptive posterior over the hidden environment state, and a clean separation between the two. Recovery of this factorisation would simultaneously validate both halves of the hypothesis and tie our result to the broader Simplex line of work on internal factorisation of independent generative processes.

This setup has several advantages for the purpose of testing our hypothesis. First, the notion of an "action" needs to be learned by the model: it is not informed ahead of time that the $a_i$ are causally downstream of its outputs. Second, we can vary the training regime to see whether the internal representations change under the meta-learning pressures we discussed above. We organise our training regimes into three conceptual groups along axes the hypothesis identifies: whether the model produces its own actions during training (passive vs active), and whether the actor at action positions is the model alone or one of a family (single-actor vs multi-actor).

\begin{enumerate}
    \item \emph{Passive baseline (Regime 1).} We train the model on token sequences emitted by the mathematically optimal Bayes-adaptive policy, which can be computed analytically as in \citep{shai2024belief}. The action tokens are external as they are not sampled from the model's own outputs. This serves as a control: the hypothesis predicts no Self representation, since the model has no causal role in producing the action stream.
    \item \emph{Active single-actor regimes (Regimes 2--4).} The model produces actions autoregressively, so the action stream is causally produced by it. We test three training-signal variants. \emph{Regime 2} applies cross-entropy on the (action, observation) stream, with two sub-regimes: (2a) places the loss at both action and observation positions, while (2b) places it only at observation positions, training the model purely as a world model conditioned on its own action history. The contrast asks whether the causal-sampling structure of self-produced tokens alone induces a Self representation, or whether explicit gradient at action positions is required. \emph{Regime 3} substitutes reinforcement learning on trajectory reward, the LLM RL-post-training analog. \emph{Regime 4} allows the model to take a budget of $K$ actions and observations before being scored on the next observation (the chain-of-thought analog), with $K$ varied from $1$ upward to probe how the exploration budget shapes representations.
    \item \emph{Mixed actors (Regime 5).} The active regimes so far always have the model itself as the actor at every action position. We now expose the model to action streams produced by a \emph{family} of policies, one of which is itself. Each trajectory is generated by first sampling an actor identity (the model itself, a random policy, the optimal Bayes-adaptive policy, an earlier checkpoint, or another deliberately distinct policy) and then having that actor produce all action tokens. The model is not told which actor produced any given trajectory and is trained on observation-prediction loss as in (2b). To predict observations well, the model must infer something about the actor's policy from the action statistics, and the natural feature is an \emph{actor-identity} representation that distinguishes the model's own action stream from those of other actors. This is to be contrasted with the \emph{token-type} self-representation (action versus observation positions) that previous regimes already encourage. This regime is also the toy analogue of the LLM training story, where the model is a predictor over many voices in pretraining and becomes the privileged producer of one of them in post-training. Crucially, by varying the actor while holding the World fixed, Regime 5 is the cleanest test of the causal form of the second hypothesis prediction: it forces the model to learn a World invariant under actor identity --- the structural signature of a genuine causal separation between self and world, distinct from the action-position / observation-position asymmetry which the other regimes already provide.
    \item \emph{Epistemic queries (Regime 6).} The actions so far are pragmatic: each one transitions the World state. We add a separate, purely epistemic channel. The model may spend a budget of query tokens, each of which returns the value of a fixed coarse-graining (epimorphic image) of the current World state without disturbing it; spending the budget well requires the model to consult the quotient it is most uncertain about. This regime targets a notion of self distinct from the previous two --- not self-as-cause or self-as-actor, but an \emph{epistemic self}: a representation of a property of the model's own belief state. The motivating question is whether choosing queries optimally forces genuine introspective access to that belief, or whether it can be done by a first-order world model alone; \texttt{strand\_three} below shows the behavioural form is deflationary but the mechanistic form is not, and lays out the dissociation and intervention tests that separate the two. This regime reuses the Regime 2b training setup, keeps full analytic ground truth, and is the toy isolation of the metacognitive question that the chain-of-thought extension of the previous strand cannot pose cleanly.
\end{enumerate}

\section*{Analysis of Trained Networks}

For each of the five training regimes above we obtain a trained transformer whose internal computations we can interrogate. Because the SWITCH is small, the exact Bayes-adaptive posterior over its World hidden state at every timestep is analytically computable; we therefore have a ground-truth comparator that experiments on natural-language reasoning models do not have. The analyses below address the two questions raised in the contribution paragraph above. \emph{First, what mechanistic representations does the model develop?} We ask whether it represents its beliefs over the World, whether it represents itself (both as a token-type position and, in Regime 5, as an actor identity), and whether the two factorise as the hypothesis predicts. \emph{Second, does the incentive structure of training causally induce these representations, and are the representations themselves causally load-bearing?} We test this by varying the training regime systematically and by performing targeted interventions on the representations once found. The first three paragraphs below address the first question and the last two address the second.

\paragraph{Belief geometry.} Following \citep{shai2024belief}, we propose training linear probes from residual-stream activations at each layer and position to the exact Bayes-adaptive posterior over the World hidden state given history. The first question is whether the mixed-state-presentation geometry recovered by \citep{shai2024belief} in the passive case generalises to the active case at all. The second is whether the posterior the model encodes is the right one: in the active regimes the Bayes-adaptive posterior must condition on the action that was taken, which is itself a sample from the model's own output distribution. If the model carries this out correctly, this would be a remarkable instance of reflexive self-knowledge in artificial agents, and a strong confirmation of the first half of the hypothesis.

% =====================================================================
% Strand Two: the geometry of the internal Self/World representations.
% Intended to slot into the "Analysis of Trained Networks" section,
% expanding the "Belief geometry" and "Self-recognition" paragraphs.
% Reuses only citation keys already present in references.bib.
% =====================================================================

\subsection*{The geometry of the Self/World factorisation}

The analyses above ask \emph{whether} the SWITCH factorisation is recovered; here we
make precise \emph{what geometric object} we expect to find, so that the probes and
interventions have a concrete target. In the passive predictor of
\citep{shai2024belief} the object is well understood: a belief is a point in the
probability simplex $\Delta_W$ over the World hidden states, the Bayes-optimal belief
given a history is a single point in $\Delta_W$, and the set of all reachable beliefs
is the (typically fractal) mixed-state presentation (MSP), the attractor of the
iterated function system whose maps are the observation-conditioned Bayes updates. Two
computations sit on top of this geometry: a \emph{linear} emission read-out giving the
next-token distribution, and a \emph{nonlinear} filtering update that moves the belief
point as each observation arrives, both of which \citep{shai2024belief} recover
linearly from the residual stream.

\paragraph{From one simplex to a product.} The SWITCH hidden state is the pair $(s,w)$,
so the belief lives in the joint simplex over $S\times W$. The central hypothesis,
stated geometrically, is that the reachable-belief manifold collapses onto the
\emph{product} submanifold $\Delta_S\times\Delta_W$ --- that the joint belief factors as
$b(s,w)\approx b_S(s)\,b_W(w)$ with the two marginals independently decodable --- rather
than filling an entangled subset of the joint simplex. This is the precise sense in
which Self and World are ``separable'', and it ties our result to the Simplex line of
work on the representation of independent generative processes. Token parity now selects
between two semantically distinct updates rather than one: an \emph{observation} token
triggers a filtering update that contracts $b_W$ toward consistency with what was seen,
while an \emph{action} token triggers a \emph{control} pushforward
$b_W \mapsto P(\cdot\mid\cdot,a)^{\top} b_W$ that advances the World belief through the
action-conditioned dynamics \citep{bengio1995iohmm}. Recovering two distinct update
maps, selected by parity and acting on a product geometry, is the first concrete target
of this strand.

\paragraph{The Self simplex as a persona-selection model.} In Regimes 2--4 the model is
always the actor, so $b_S$ is a point mass: there is no uncertainty about who is acting
and the Self marginal is degenerate, leaving only the thin token-type (action vs.\
observation) distinction. The Self simplex becomes a non-trivial object precisely in the
mixed-actor Regime 5, where each trajectory is generated by one of a family of policies
$\{\pi^{(1)},\dots,\pi^{(M)}\}$, one of which ($\pi^{(0)}$) is the model itself. Here
$\Delta_S$ is the online posterior over \emph{actor identity}: its vertices are the
members of the actor family, and it is updated by the likelihood ratio of each observed
action under the candidate policies --- exactly a filtering process on a ``which-actor''
latent. We therefore expect a second, lower-dimensional mixed-state geometry --- a
persona-selection MSP --- to appear alongside the World MSP, linearly decodable from the
residual stream in Regime 5 but absent in Regimes 2--4.

\paragraph{Nested belief and theory of mind.} The persona simplex does more than label
trajectories: it gates the World update. An action is evidence about the World state
only insofar as the acting policy is World-state-dependent, and how strongly depends on
which actor produced it. The action-token update on $b_W$ therefore carries an
actor-dependent Bayesian factor \emph{before} the pushforward,
\[
  b_W(w)\;\propto\; b_W(w)\,\pi^{(c)}(a\mid w)\quad\text{(actor-gated)},
  \qquad\text{then}\qquad
  b_W \;\mapsto\; P(\cdot\mid\cdot,a)^{\top} b_W,
\]
so that a random-actor action is treated as uninformative (pushforward only) while an
optimal-actor action drives a strong belief update. When the actors are themselves
\emph{stateful} --- e.g.\ Bayes-adaptive agents carrying their own World belief ---
predicting their actions requires the model to simulate \emph{that actor's} belief, a
separate point in a belief simplex that need not coincide with the model's own. The full
latent geometry is then a \emph{bundle}: over each non-self node $c$ of $\Delta_S$ sits a
fiber simplex $\Delta_W^{(c)}$ encoding the model's representation of actor $c$'s
beliefs, distinct both from the ground-truth World state and from the model's own
belief. Representing another agent's belief as separable from one's own is first-order
theory of mind; the SWITCH with stateful actors forces it as a computational requirement
of observation prediction, and it is jointly testable with the World geometry on the
same analytic ground truth.

\paragraph{The self as a degenerate fiber.} The self node is distinguished not by its
location in $\Delta_S$ but by the status of its fiber: over the self node the belief
fiber is \emph{identified} with the model's own World belief $b_W$ rather than separately
inferred, because the self acts on its own belief by construction. This identification is
the geometric content of the reflexive self-recognition reported behaviourally in
post-trained models \citep{betley2025tellme,panickssery2024self} and in our forthcoming
work with Lindsey. It is also what makes the Self component causally load-bearing: a
perturbation written into the self fiber should move both the model's observation
predictions and its own action distribution, whereas a perturbation to a non-self fiber
$\Delta_W^{(c)}$ should move only the predicted behaviour of actor $c$. This double
dissociation, layered on top of the Self-steering interventions already described, is the
sharpest causal signature of a genuine self representation.

\paragraph{Circuit-level predictions and the off-policy test.} Two implementations are
distinguishable and worth separating experimentally. Under a \emph{shared, identity-gated}
update circuit, a single mechanism performs the actor-conditioned belief update with the
$\Delta_S$ vector entering as a gating input; steering $\Delta_S$ should then continuously
deform the observed update rule, from ``ignore the action as evidence'' to ``treat it as
strong evidence''. Under a \emph{privileged self pathway}, the self is handled by the
first-person generation circuit --- read off, not inferred --- rather than by the shared
theory-of-mind module, and only the non-self nodes route through the latter.
Distinguishing these is the mechanistic content of the off-policy generalisation test
already proposed: a shared identity-gated bundle should interpolate to a held-out actor
of a known type (Hypothesis A, generalisation), whereas per-actor shortcuts through the
joint (action, observation) statistics should not (Hypothesis B). Two design notes follow
directly: the nested fiber structure arises only if the actor family contains stateful
policies (memoryless actors require actor identity but no fiber), and the actors should
remain first-order World modellers in the toy --- recursive actor-models-actor structure
is the layered self-simulation we defer to the language-model extension.


% =====================================================================
% Strand Three: epistemic actions and introspective access.
% Regime 6: the model may query lumped quotients (epimorphic images)
% of the World HMM. The analysis separates genuine introspective
% access to the model's own belief state from first-order competence
% that merely uses it.
% =====================================================================

\subsection*{Epistemic actions: queryable quotients and introspective access}

The regimes above vary \emph{who} acts; this strand varies \emph{what acting is
for}. Following the classical distinction of \citep{kirsh1994epistemic}, the
SWITCH actions of Regimes 2--5 are simultaneously \emph{pragmatic} (they
transition the World state) and \emph{epistemic} (the response is informative
about it). Regime 6 adds a channel that is purely epistemic: it yields
information about $w$ without disturbing it.

\paragraph{Setup (Regime 6).} Fix a family of epimorphisms
$\varphi_i : W \to W^{(i)}$ onto smaller HMMs, chosen so that each quotient is a
lumpable coarse-graining of the World chain \citep{kemeny1960finite}. At
designated positions the model may emit a query token $q_i$; the environment
answers with the current coarse state $y = \varphi_i(w_t)$, leaving $w_t$
untouched. Queries are budgeted (a fixed number per trajectory, or a small
reward cost), so the model must choose \emph{which} quotient to consult. The
cleanest instantiation factors the World outright, $W \cong A \times B$, with
the $\varphi_i$ the two projections: the model holding belief $b_W$ must decide
whether it is more ignorant about the $A$ factor or the $B$ factor. The
intuition motivating the regime is that spending the budget well requires the
model to \emph{know what it knows}: a metacognitive judgement about its own
belief state, made before any new evidence arrives.

\paragraph{The deflationary observation.} We state at the outset the objection
that this intuition does not survive unmodified, because the objection shapes
the entire analysis plan. For an ideal Bayes-adaptive agent the expected
information gain of query $i$ is the entropy of the pushforward of its current
belief, $\mathrm{EIG}_i = H(\varphi_{i\,\sharp}\, b_W)$: the answer is a
deterministic function of $w$, so the mutual information $I(w;\varphi_i(w))$
reduces to the entropy of the coarse marginal. The optimal query policy
$\arg\max_i H(\varphi_{i\,\sharp}\, b_W)$ is therefore a functional of the
first-order belief --- one more readout of $b_W$, requiring no representation
\emph{of} the belief over and above the belief itself. Optimal query behaviour
alone is thus compatible with a pure world model. This is precisely the
deflationary critique long aimed at animal and infant metacognition: opt-out
and help-seeking paradigms \citep{hampton2001monkeys,goupil2016infants} admit
first-order reconstructions in which an uncertainty signal is \emph{used}
without being \emph{represented} \citep{carruthers2008skeptical}, and
behavioural evidence alone has not resolved that dispute. The toy setting
cannot evade the equivalence at the behavioural level --- but it can break it
at the mechanistic level, in two ways.

\paragraph{Escape one: dissociation under miscalibration.} The behavioural
equivalence holds only for the ideal agent, for whom ``my uncertainty'' is a
deterministic function of the observable history. A finite trained model's
realised belief $\hat b_W$ --- the object our probes decode --- deviates from
the analytic posterior $b_W$ somewhere: off the training distribution, on long
or ambiguous histories, at capacity limits. There the hypotheses come apart.
\emph{Introspective access} predicts queries that track the model's own
miscalibration, $\arg\max_i H(\varphi_{i\,\sharp}\, \hat b_W)$: the model asks
about what \emph{it} is uncertain of, even where an ideal observer would not
be. A \emph{compiled heuristic} --- a query policy fit to the typical
uncertainty profile of surface history features --- instead approximates the
ideal policy on-distribution and tracks neither posterior where they diverge.
Because both $b_W$ (analytically) and $\hat b_W$ (by the belief-geometry
probes) are available to us, we can locate histories on which they disagree and
classify every query the model makes. Knowing-what-you-know, on this
operationalisation, is not being right about the world; it is being right about
one's own errors.

\paragraph{Escape two: causal introspection.} The sharper test is
interventional and reuses machinery this proposal already builds. Once the
World-belief subspace is identified, we edit $\hat b_W$ in the residual stream
--- writing in a belief with a different marginal-entropy profile --- and ask
whether the query distribution shifts as the edited belief predicts. A positive
result establishes a causal pathway from the belief representation to the query
choice: the metacognitive act \emph{reads} the model's own internal state
rather than recomputing a shortcut from the history. This is the toy-model
analogue of concept-injection tests of introspection in language models, and
connects directly to the on-policy-detection findings in our forthcoming work
with Lindsey. We note the regress this stops: one can always redescribe
``reading one's own belief state'' as more first-order computation, but at that
point the deflation has no further content --- a circuit that consumes the
belief representation \emph{as} a representation, tracks its idiosyncratic
errors, and is causally load-bearing for behaviour is everything introspective
access could be in any physical system.

\paragraph{The epistemic self and its geometry.} If the introspective readout
exists, Regime 6 adds a third notion of self to the two already in play
(token-type and actor-identity): an \emph{epistemic self}, a represented
quantity whose content is a property of the model's own belief state rather
than of the World. In the factored instantiation the predicted object is
concrete: scalar readouts of the marginal entropies $H(b_A)$ and $H(b_B)$,
linearly available upstream of the query decision. These quantities sit on the
Self side of the Markov blanket while being \emph{about} the World --- the
first genuinely second-order representation the SWITCH demands, and the
self-directed counterpart of the other-directed belief fibers of the previous
strand. This symmetry raises a sharp circuit-level question, parallel to the
privileged-self-pathway question raised there: does the self-uncertainty
readout route through the same machinery that (with stateful actors) represents
\emph{other} agents' beliefs, or through a privileged direct pathway? The
former is the mechanistic form of the thesis that metacognition is self-directed
mindreading \citep{carruthers2009mindreading}; the SWITCH with both stateful
actors and queryable quotients makes that long-standing question empirical.

\paragraph{Design notes.} (i) Query timing is curriculum-staged: fixed query
slots first, so the model chooses only \emph{which} quotient to consult, with
learned timing and explicit query costs as a second phase. This keeps the
regime free of the cold-start problem that makes open-ended reasoning training
expensive. (ii) The quotient family should be designed so that the identity of
the most informative quotient varies across histories (the product projections
already guarantee this, as uncertainty migrates between factors); otherwise a
static preference over oracles is optimal and the regime tests nothing.
(iii) The regime reuses the Regime 2b training infrastructure, and its analytic
comparators --- filtered posterior, pushforward entropies, value of information
--- are cheap to compute. (iv) A natural extension, deferred on cost grounds,
replaces the external oracle with free-form internal rollouts in which the
model generates both halves of a simulated exchange (the chain-of-thought
analogue); the quotient-oracle version isolates the metacognitive question
without first having to train reasoning from scratch.


\paragraph{Self-recognition and off-policy generalisation.} Actor-identity self-recognition is not predicted in Regimes 2--4, where the model is always the actor and has no contrastive training signal to distinguish itself from another agent. Regime 5 changes this by exposing the model to action streams from many policies including its own, giving the actor-identity Self representation a functional purpose. The first test, then, is whether ``this trajectory was produced by my own policy'' is linearly decodable from the residual stream in Regime 5 but not in Regimes 2--4. A positive result in Regime 5 is the toy-model analogue of the on-policy detection finding in our forthcoming work with Lindsey, and is the expected outcome if the actor-identity Self has formed. The harder test is off-policy generalisation: when presented with an action stream from a policy not in the training family, does the model produce correct Bayes-adaptive predictions? This distinguishes two hypotheses about the kind of computation the model has learned. Under Hypothesis A, the model implements the abstract action-conditioned belief update, so the source of the action does not matter and generalisation follows. Under Hypothesis B, the model has learned a policy-specific shortcut through learning the joint statistics of (action, observation) under the policies it has seen and generalises poorly.

\paragraph{Cross-world generalisation.} If the Self and World components are genuinely factored, they should be modifiable independently. We test this by training the model in a meta-learning regime over a family of SWITCHes sharing the same World-state cardinality and observation alphabet, but with transition matrices sampled from a meta-distribution and presenting a held-out SWITCH from the same family at evaluation. The hypothesis predicts that the World representation adapts in-context to the new task while the Self representation remains stable. This is a toy analogue of learning the self as a fixed actor in a volatile world.

\paragraph{Cross-regime training dynamics.} The probes above can be applied not only to fully-trained models but to checkpoints throughout training, giving us the time-course of representation emergence in each regime. The hypothesis predicts a specific dissociation pattern. World-state belief should sharpen in all five regimes, since it is required for any observation prediction. The two notions of Self representation we distinguished above should dissociate across regimes. A \emph{token-type} Self, distinguishing action positions from observation positions for the purposes of belief updates, should emerge in Regimes 2--4, where the model produces its own actions; the action-position entropy collapse should be sharpest in (2a). An \emph{actor-identity} Self, distinguishing the model's own action stream from those of other actors, should emerge specifically in Regime 5, where the model is exposed to multiple actors and must distinguish among them to predict observations. A clean sub-question within token-type Self is the contrast between (2a) and (2b): if a Self representation emerges in (2b), where the loss never directly touches action positions, then the causal-sampling structure of the input is itself sufficient; if only in (2a), then explicit gradient at action positions is required. Either outcome refines what we should expect to find in RL-trained reasoning LLMs, where both pressures are typically present at once.

\paragraph{Causal control via Self-steering.} To establish that the SWITCH factorisation is causal, we perform targeted interventions on the Self and World subspaces identified above. Writing a different Self encoding into the residual stream should shift the action distribution in the direction the SWITCH dynamics predict, without disturbing the World belief; writing a perturbed World encoding should distort observation predictions and downstream action choices in the direction the Bayes-adaptive update predicts. A clean dissociation across these interventions is strong evidence that the SWITCH factorisation is real and not merely linearly decodable, and provides a toy-model existence proof that the Self component of an analogous LLM representation could be a steering handle for interpretability and alignment.

\section*{Extension to Language Models (Second Half of the Year)}

The second half of the grant is conditional on the first half delivering interpretable findings on the SWITCH. Assuming it does, we extend the probe suite to RL-trained reasoning language models, but in a setting that preserves the genuine self/world interaction structure the hypothesis depends on. Pure chain-of-thought reasoning is not the best setting for such an experiment: the model's own tokens are its only input, and there is no external World with hidden state and action-conditioned dynamics for the model to attribute to. We focus instead on \emph{code-execution-augmented reasoning}, where the model interleaves natural-language reasoning with code that is actually executed and whose output is fed back into the context. The Python interpreter plays the role of World, with its own hidden state and action-conditioned responses, and the model's code blocks play the role of Self-emitted actions. Several open-source checkpoints exist in this configuration with a paired base model available (e.g., from the rStar-Math / Qwen-Math-with-code family), and we can check for the emergence of self/world factorizations during post-training. The SWITCH framework here serves not only as a target structure to look for, but as a scaffold for asking richer questions that the toy version cannot pose: whether the Self subspace shifts character between exploratory and committed modes of reasoning (e.g., during chain-of-thought versus when committing to a final answer), whether the model handles layered self-simulation --- role-playing another agent while remaining itself --- as a distinguishable form of self separate from the binary self / not-self of Regime 5, and whether the model's trust in tool output is calibrated against its own beliefs, so that anomalous or fabricated tool responses are integrated with appropriate scepticism rather than as ground truth. These are LLM-specific phenomena that the toy SWITCH cannot reach, but for which the toy work supplies the probes that let us ask after them precisely.

Taken together, the proposed project offers a theoretical framework in the SWITCH model for a Self/World factorisation in intelligent agents, a mechanistic hypothesis for how this factorisation emerges from training pressures, and a concrete experimental design to test the hypothesis in a toy model and then in language models. This builds on a conceptual distinction between passive prediction and active exploration that is central to agentic learning, and a recognition of the special role that exploration plays in shaping the internal representations of an agent, and especially its self model. The project also has implications for interpretability and alignment, since the self-representation we predict is a natural handle for understanding and influencing the agent's behaviour. Last but not least, the project's results will have implications for developmental cognitive science, since they will clarify the computational mechanisms that can produce the patterns of behaviour observed in infants and foetuses, and sharpen the hypotheses that developmental researchers are themselves testing. 

\bibliography{references}

\end{document}
